%==============================================================================
% GUION DE SUSTENTACIÓN — 45 MINUTOS, DOS PONENTES
% Tratamiento Numérico de la Inestabilidad de Kelvin–Helmholtz
% bajo el Marco de la Magnetohidrodinámica Relativista Resistiva
%
% Compilación: pdflatex guion_sustentacion.tex  (×2)
% El guion sigue lámina a lámina el deck presentacion/main.pdf (50 pp,
% 31 láminas de cuerpo + referencias + agradecimiento + respaldo B1–B13).
%==============================================================================
\documentclass[11pt,letterpaper]{article}

\usepackage[utf8]{inputenc}
\usepackage[T1]{fontenc}
\usepackage[spanish,es-tabla]{babel}
\usepackage{lmodern}
\usepackage[margin=2.4cm]{geometry}
\usepackage{xcolor}
\usepackage{booktabs}
\usepackage{array}
\usepackage{enumitem}
\usepackage[most]{tcolorbox}
\usepackage{fancyhdr}
\usepackage{lastpage}
\usepackage[hidelinks]{hyperref}

%---------------------------------------------------------- paleta (coherente con el deck)
\definecolor{primary}{HTML}{1B0C41}   % púrpura Inferno profundo
\definecolor{accent}{HTML}{E65D2F}    % naranja Inferno
\definecolor{gold}{HTML}{F5B841}
\definecolor{soft}{HTML}{F7F1E1}
\definecolor{grisnota}{HTML}{555555}

%---------------------------------------------------------- encabezado / pie
\pagestyle{fancy}
\fancyhf{}
\fancyhead[L]{\footnotesize\itshape Guion de sustentación --- KHI en RRMHD}
\fancyhead[R]{\footnotesize Martinez \& Rodriguez}
\fancyfoot[C]{\footnotesize Página \thepage\ de \pageref{LastPage}}
\renewcommand{\headrulewidth}{0.4pt}

%---------------------------------------------------------- cajas
% Parlamento: el texto que se pronuncia literalmente.
\newtcolorbox{parlamento}[1]{%
  breakable, enhanced,
  colback=soft!35, colframe=primary,
  boxrule=0.6pt, arc=1.5pt,
  left=8pt, right=8pt, top=6pt, bottom=6pt,
  title={\small\bfseries #1 --- texto para pronunciar},
  fonttitle=\color{white},
  colbacktitle=primary}

% Indicaciones escénicas y técnicas (no se pronuncian).
\newtcolorbox{indicaciones}{%
  breakable,
  colback=white, colframe=accent,
  boxrule=0.5pt, arc=1.5pt,
  left=8pt, right=8pt, top=4pt, bottom=4pt,
  title={\small\bfseries Indicaciones (no se pronuncian)},
  fonttitle=\color{white},
  colbacktitle=accent}

% Encabezado de lámina
\newcommand{\lamina}[5]{%
  \subsection*{Lámina #1 --- #2}
  \noindent
  \begin{tabular}{@{}llll@{}}
    \textbf{Ponente:} #3 &\quad \textbf{Duración:} #4 &\quad
    \textbf{Tiempo acumulado al cierre:} #5 &
  \end{tabular}
  \vspace{2pt}}

\newcommand{\relevo}[1]{%
  \begin{center}
  \begin{tcolorbox}[colback=gold!20,colframe=gold!80!black,boxrule=0.6pt,
    width=0.92\linewidth,arc=1.5pt,left=6pt,right=6pt,top=4pt,bottom=4pt]
    \small\textbf{Relevo de ponente.} #1
  \end{tcolorbox}
  \end{center}}

\setlength{\parskip}{4pt}
\setlength{\parindent}{0pt}

%==============================================================================
\begin{document}

%---------------------------------------------------------- portada
\begin{titlepage}
\centering
\vspace*{1.2cm}
{\color{primary}\rule{\linewidth}{2pt}}\\[0.9cm]
{\large Universidad Distrital Francisco José de Caldas}\\[2pt]
{\small Facultad de Ciencias Matemáticas y Naturales --- Programa Académico de Física}\\[1.6cm]
{\LARGE\bfseries Guion de Sustentación}\\[0.5cm]
{\large Trabajo de Grado}\\[0.9cm]
{\color{primary}\Large\bfseries Tratamiento Numérico de la Inestabilidad de
Kelvin--Helmholtz bajo el Marco de la Magnetohidrodinámica Relativista
Resistiva}\\[1.4cm]
\begin{tabular}{rl}
\textbf{Autores y ponentes:} & Bryan Martinez Anzola\\
                             & Sebastián Rodriguez Garcia\\[6pt]
\textbf{Director:}           & Prof.\ Ph.D.\ Sergio Miranda-Aranguren\\[6pt]
\textbf{Jurado:}             & Prof.\ Juan Carlos Giraldo Acuña\\[6pt]
\textbf{Duración prevista:}  & 45 minutos (exposición) $+$ turno de preguntas\\[6pt]
\textbf{Material:}           & \texttt{presentacion/main.pdf} (50 páginas;
                               31 láminas de cuerpo,\\
                             & referencias y material de respaldo B1--B13)\\
\end{tabular}\\[1.6cm]
{\color{primary}\rule{\linewidth}{2pt}}\\[0.8cm]
{\small Bogotá, D.C.\ --- Julio de 2026}
\vfill
\end{titlepage}

\tableofcontents
\newpage

%==============================================================================
\section{Instrucciones generales de uso}
%==============================================================================

\subsection{Naturaleza del documento}
El presente guion acompaña, lámina por lámina, la presentación
\texttt{main.pdf} de la sustentación. Para cada lámina se indica el ponente a
cargo, la duración asignada, el tiempo acumulado al cierre, el \emph{texto
para pronunciar} (registro formal, redactado para ser leído o memorizado con
fidelidad conceptual, no necesariamente literal) y las \emph{indicaciones}
escénicas o técnicas que no se pronuncian. El texto está calibrado a un ritmo
de dicción de 130--150 palabras por minuto, apropiado para una exposición
académica formal.

\subsection{Distribución de ponentes}
Se alterna por bloques temáticos, de modo que ambos autores demuestren dominio
de la totalidad del trabajo y el relevo coincida siempre con una transición
natural del argumento. La asignación propuesta es simétrica y puede
intercambiarse de común acuerdo sin alterar tiempos ni transiciones.

\begin{center}
\begin{tabular}{@{}clccc@{}}
\toprule
Segmento & Bloque temático & Láminas & Ponente & Duración\\
\midrule
I   & Apertura, motivación y diseño          & 1--5   & Bryan Martinez & 7:30\\
II  & Marco teórico RRMHD                    & 6--11  & Sebastián Rodriguez      & 9:30\\
III & Implementación numérica y montaje      & 12--17 & Bryan Martinez & 7:30\\
IV  & Resultados I: barrido en conductividad & 18--26 & Sebastián Rodriguez      & 13:00\\
V   & Resultados II: campaña magnética; cierre & 27--32 & Bryan Martinez & 7:00\\
\midrule
\multicolumn{4}{r}{Total de exposición} & 44:30\\
\multicolumn{4}{r}{Margen de contingencia} & 0:30\\
\bottomrule
\end{tabular}
\end{center}

Totales por ponente: Bryan Martinez, 22:00; Sebastián Rodriguez, 22:30.

\subsection{Puntos de control de ensayo}
Al cronometrar los ensayos deben verificarse los siguientes hitos. Si en un
punto de control se acumula un retraso superior a un minuto, se comprimen las
láminas marcadas como \emph{comprimibles} (láminas 15, 19 y 22), que admiten
una versión de treinta segundos sin pérdida del argumento central.

\begin{center}
\begin{tabular}{@{}lc@{}}
\toprule
Hito & Minuto objetivo\\
\midrule
Fin de motivación y diseño (lámina 5)            & 7:30\\
Fin del marco teórico (lámina 11)                & 17:00\\
Fin de implementación y montaje (lámina 17)      & 24:30\\
Fin de resultados del barrido (lámina 26)        & 37:30\\
Fin de la campaña magnética (lámina 29)          & 41:30\\
Cierre de la exposición (lámina 32)              & 44:30\\
\bottomrule
\end{tabular}
\end{center}

\subsection{Protocolo de relevos}
Cada relevo se ejecuta con una frase puente pronunciada por el ponente
saliente, que cede la palabra nombrando el bloque siguiente; el ponente
entrante inicia sin saludo adicional. Los cuatro relevos están escritos en
el cuerpo del guion en recuadros dorados. Durante el segmento del compañero,
el ponente en espera permanece atento al cronómetro y opera, si así se
acuerda, el avance de láminas.

\subsection{Convenciones de notación al hablar}
En todo el guion, y en coherencia con la monografía: $W$ denota el factor de
Lorentz; $\Gamma$, exclusivamente el índice adiabático; $\hat\omega$, la tasa
de crecimiento normalizada; y $\sigma$ se reserva de manera exclusiva para la
conductividad eléctrica. El código se nombra «Cueva», sin artículo. No se
emplea jerga interna del repositorio.

%==============================================================================
\section{Segmento I --- Apertura, motivación y diseño (Bryan Martinez, 7:30)}
%==============================================================================

\lamina{1}{Portada}{Bryan Martinez}{1:00}{1:00}

\begin{parlamento}{Bryan Martinez}
Muy buenos días. Señor jurado, profesor Juan Carlos Giraldo; señor director,
profesor Sergio Miranda-Aranguren; señoras y señores. Mi nombre es Bryan
Martinez Anzola y me acompaña mi compañero Sebastián Rodriguez Garcia. Presentamos
a ustedes la sustentación de nuestro trabajo de grado, titulado
\emph{Tratamiento Numérico de la Inestabilidad de Kelvin--Helmholtz bajo el
Marco de la Magnetohidrodinámica Relativista Resistiva}, dirigido por el
profesor Miranda-Aranguren.

En una frase, este trabajo cuantifica cómo la resistividad finita y la
geometría del campo magnético gobiernan la inestabilidad de
Kelvin--Helmholtz relativista, mediante más de sesenta simulaciones
numéricas realizadas con el código Cueva.
\end{parlamento}

\begin{indicaciones}
Postura frontal, sin leer la lámina. La frase-tesis final se pronuncia con
pausa previa: es el enunciado que el jurado debe retener. Verificar de
antemano proyector, apuntador y cronómetro visible.
\end{indicaciones}

\lamina{2}{Estructura de la Presentación}{Bryan Martinez}{1:30}{2:30}

\begin{parlamento}{Bryan Martinez}
La exposición sigue la estructura que ustedes observan. Primero, la
motivación física y astrofísica del problema; segundo, el marco teórico de la
magnetohidrodinámica relativista resistiva; tercero, la implementación
numérica y el montaje experimental; y finalmente los resultados, organizados
en dos partes, con sus conclusiones.

Quiero destacar desde ahora los cuatro objetivos específicos aprobados en la
propuesta, porque los resultados y las conclusiones se presentarán en
correspondencia uno a uno con ellos. Primero, cuantificar el efecto de la
resistividad sobre la tasa de crecimiento lineal de la inestabilidad.
Segundo, analizar su efecto en la formación de estructuras secundarias:
vórtices y láminas de corriente. Tercero, evaluar el impacto de la difusión
magnética sobre las variables electromagnéticas globales. Y cuarto, evaluar
cómo la intensidad y la orientación del campo magnético externo modifican la
evolución de la inestabilidad. La herramienta es el código Cueva, un código
de magnetohidrodinámica relativista resistiva en volúmenes finitos, y la base
de datos comprende cuarenta y cuatro simulaciones del barrido en
conductividad más dieciocho configuraciones magnéticas.
\end{parlamento}

\begin{indicaciones}
Señalar la caja de objetivos al enumerarlos, marcando con la mano uno a
cuatro. La correspondencia objetivo--resultado--conclusión es el hilo
argumentativo que el jurado debe ver desde el primer minuto.
\end{indicaciones}

\lamina{3}{El Fenómeno: la KHI en el Universo Relativista}{Bryan Martinez}{2:00}{4:30}

\begin{parlamento}{Bryan Martinez}
Comencemos por el fenómeno, sin ecuaciones. La inestabilidad de
Kelvin--Helmholtz aparece dondequiera que dos flujos en contacto se mueven a
velocidades distintas. Esa capa de cizalla constituye un reservorio de
energía libre. Si la interfaz se ondula levemente, el flujo se acelera sobre
las crestas y se frena en los valles; por el principio de Bernoulli, esa
diferencia de velocidades genera gradientes de presión que ejercen un torque
sobre la interfaz. El torque amplifica la ondulación, la ondulación
intensifica el torque, y el resultado es un crecimiento exponencial: la fase
lineal. Cuando la amplitud deja de ser pequeña, la interfaz se enrolla en los
vórtices característicos que ustedes observan en la figura: la fase no
lineal.

En el universo relativista este mecanismo es ubicuo: opera en las fronteras
de los chorros de núcleos galácticos activos y de los estallidos de rayos
gamma, donde un flujo cercano a la velocidad de la luz roza un medio en
reposo. Allí gobierna la mezcla y el frenado del chorro ---el llamado
\emph{entrainment}, vinculado a la dicotomía morfológica entre radiofuentes
FR-I y FR-II--- y actúa además como dínamo: los vórtices estiran las líneas
de campo magnético y lo amplifican. Entender qué controla esta inestabilidad
es, por tanto, entender un regulador central de la dinámica de los chorros
astrofísicos.
\end{parlamento}

\begin{indicaciones}
Esta lámina es, a la vez, la respuesta ensayada a la pregunta pedagógica
«explíquela sin fórmulas» (véase el Anexo A.1). Señalar la figura del chorro
al mencionar las fronteras de los jets.
\end{indicaciones}

\lamina{4}{Por Qué Resistiva: los Límites del Congelamiento Ideal}{Bryan Martinez}{1:30}{6:00}

\begin{parlamento}{Bryan Martinez}
Ahora bien, ¿por qué resistiva? La ecuación de inducción contiene dos
términos: advección y difusión, y la difusión está pesada por la
resistividad, que es el inverso de la conductividad. En el límite de
conductividad infinita ---la magnetohidrodinámica ideal--- rige el teorema de
congelamiento de Alfvén: la topología del campo magnético no puede cambiar y
la reconexión está prohibida. Sin embargo, cuando un código ideal muestra
reconexión, esta proviene del error de truncamiento de la malla: es una
resistividad numérica, incontrolada, que depende de la resolución y no de la
física.

El marco resistivo invierte esa situación: al introducir una conductividad
finita como parámetro físico, la reconexión y el calentamiento óhmico se
vuelven procesos naturales, predictivos y controlados, imprescindibles en
escenarios como las llamaradas de núcleos activos o las magnetosferas de
púlsares. De aquí nace la pregunta central de este trabajo: ¿cómo gobiernan
la conductividad finita y la geometría del campo magnético el crecimiento, la
morfología y el balance energético de la inestabilidad de Kelvin--Helmholtz
relativista?
\end{parlamento}

\begin{indicaciones}
Aquí se planta, desde la motivación, la respuesta a la objeción «¿su
disipación es física o numérica?»; la evidencia empírica se anunciará en la
lámina 26. Enfatizar el contraste «incontrolado» frente a «controlado».
\end{indicaciones}

\lamina{5}{Diseño del Experimento Numérico: Dos Campañas}{Bryan Martinez}{1:30}{7:30}

\begin{parlamento}{Bryan Martinez}
Para responder esa pregunta diseñamos dos campañas sobre un único montaje
base. La primera campaña aborda los objetivos uno a tres: un barrido de
cuarenta y cuatro simulaciones en las que la conductividad recorre el
intervalo de cien a diez mil quinientos, con todos los demás parámetros
fijos; de este modo se aísla el efecto de la difusividad magnética sobre la
tasa lineal, la fenomenología no lineal y los diagnósticos
electromagnéticos globales.

La segunda campaña aborda el objetivo cuatro: se fija la conductividad en dos
regímenes ---seis mil y diez mil--- y se varía el campo magnético en tres
familias. La familia A varía la intensidad; la familia B rota el campo en el
plano $y$--$z$, activando progresivamente la tensión magnética, porque
aparece componente paralela al vector de onda; y la familia C rota el campo
en el plano $x$--$z$, donde el producto escalar entre el vector de onda y el
campo es nulo y, por tanto, no hay tensión: es nuestro experimento de
control. Con esta arquitectura, cedo la palabra a mi compañero Sebastián
Rodriguez, quien presentará el marco teórico del sistema.
\end{parlamento}

\relevo{Bryan Martinez concluye el segmento I con la frase anterior y
Sebastián Rodriguez toma la palabra directamente, sin saludo adicional.
Minuto de control: 7:30.}

%==============================================================================
\section{Segmento II --- Marco teórico RRMHD (Sebastián Rodriguez, 9:30)}
%==============================================================================

\lamina{6}{Estructura Covariante del Sistema RRMHD}{Sebastián Rodriguez}{2:00}{9:30}

\begin{parlamento}{Sebastián Rodriguez}
Gracias, Bryan. Antes de escribir cualquier ecuación, fijemos las
convenciones que regirán toda la presentación: unidades de
Heaviside--Lorentz, con la velocidad de la luz y las permitividades iguales a
uno; signatura métrica menos, más, más, más; $W$ denotará el factor de
Lorentz y la letra $\Gamma$ se reserva exclusivamente para el índice
adiabático, en la formulación de Valencia.

El sistema descansa sobre cuatro leyes covariantes. La conservación del
número bariónico; la conservación del tensor de energía-momento
\emph{total}, suma del término de fluido perfecto y del término
electromagnético de Faraday--Maxwell, ambos obtenidos por variación de la
métrica según el procedimiento de Hilbert; las ecuaciones de Maxwell
inhomogéneas, que sí son dinámica, pues provienen de mínima acción con
fuente; y las homogéneas. Sobre estas últimas conviene un punto fino: dado
que el tensor de Faraday es la derivada exterior del cuadripotencial, su
derivada exterior se anula por nilpotencia. Es decir, la ausencia de
monopolos magnéticos no es una ecuación dinámica sino una identidad
geométrica, la identidad de Bianchi. Esta distinción, que parece un
tecnicismo, será decisiva cuando pasemos al problema discreto. Las
derivaciones completas están disponibles en el material de respaldo.
\end{parlamento}

\begin{indicaciones}
Las convenciones se declaran aquí una única vez y no se vuelven a mencionar.
La distinción identidad geométrica frente a ecuación dinámica prepara la
lámina 8 y responde por anticipado al examen tensorial del jurado
(respaldos B1 y B2).
\end{indicaciones}

\lamina{7}{La Ley de Ohm Relativista: el Cierre Resistivo}{Sebastián Rodriguez}{2:00}{11:30}

\begin{parlamento}{Sebastián Rodriguez}
El sistema anterior no está cerrado: falta la ecuación constitutiva que
define el régimen resistivo, la ley de Ohm relativista. En el marco comóvil,
para una respuesta lineal e isótropa, la corriente es la suma del término
convectivo de carga y del término conductivo, proporcional a la contracción
del tensor de Faraday con la cuadrivelocidad. Al proyectarla sobre la malla
euleriana en el formalismo tres más uno se obtiene la expresión enmarcada:
nótese que no es una suma galileana; el factor de Lorentz y los
acoplamientos transversales son la huella genuinamente relativista.

Esta ley contiene los dos límites físicos del problema. Cuando la
conductividad tiende a infinito, el campo eléctrico comóvil se anula y se
recupera el congelamiento ideal. Cuando la conductividad es finita, el campo
eléctrico comóvil sobrevive, y con él el calentamiento de Joule y la
reconexión. Declaramos aquí, con honestidad, los límites de validez del
cierre: la conductividad es escalar e isótropa, lo que presupone un plasma
colisional, con frecuencia de colisiones mucho mayor que la frecuencia de
ciclotrón, sin efecto Hall. La resistividad es un coeficiente de transporte
macroscópico que emerge de las integrales de colisión del nivel cinético de
Vlasov--Boltzmann. Y este cierre tiene una consecuencia dinámica que
condiciona todo el trabajo numérico: el término fuente proporcional a la
conductividad vuelve el sistema rígido cuando ella es grande. Ese es el reto
central de la sección siguiente.
\end{parlamento}

\begin{indicaciones}
La frase sobre el origen cinético de la resistividad es una respuesta
ensayada (Anexo A.4; respaldos B5 y B9). Señalar la ecuación enmarcada al
pronunciar «expresión enmarcada».
\end{indicaciones}

\lamina{8}{Restricciones de Divergencia: el Sistema Aumentado GLM}{Sebastián Rodriguez}{1:30}{13:00}

\begin{parlamento}{Sebastián Rodriguez}
Retomo la identidad de Bianchi. En el continuo, la nilpotencia de la derivada
exterior garantiza que la divergencia del campo magnético sea nula. Pero los
operadores discretos de una malla no son nilpotentes: el error de
truncamiento genera monopolos espurios que, sin tratamiento, se acumulan y
destruyen la estabilidad de la simulación.

La solución adoptada es el método de multiplicadores de Lagrange
generalizados, GLM: se introducen dos campos escalares que acoplan las
restricciones a la evolución, y puede demostrarse que el error de divergencia
obedece la ecuación del telégrafo. Esa ecuación tiene dos propiedades
exactamente deseables: es hiperbólica, de modo que el error se propaga a la
velocidad de limpieza, igual a la velocidad de la luz, sin violar jamás el
cono causal; y es disipativa, de modo que el error decae exponencialmente.
La figura ilustra el contraste: sin tratamiento el error se acumula; con GLM
se propaga fuera del dominio y se amortigua, sin alterar las soluciones
físicas. Con esto respondemos de manera constructiva cómo se garantiza la
ausencia de monopolos sin destruir la causalidad.
\end{parlamento}

\begin{indicaciones}
Esta lámina es la respuesta completa a la pregunta prevista del jurado sobre
divergencia y causalidad (Anexo A.2; respaldo B3 para el origen
variacional).
\end{indicaciones}

\lamina{9}{El Sistema en Forma Fuertemente Conservativa}{Sebastián Rodriguez}{1:00}{14:00}

\begin{parlamento}{Sebastián Rodriguez}
Todo lo anterior se ensambla en un sistema de leyes de balance con catorce
variables conservadas: densidad relativista, momento, energía, los campos
magnético y eléctrico, la carga y los dos escalares de limpieza. Subrayo una
consecuencia de conservar el tensor total: el momento y la energía incluyen
el vector de Poynting y la densidad de energía electromagnética, de modo que
la fuerza de Lorentz no aparece como término fuente, sino absorbida en la
divergencia del flujo; esto garantiza conservación exacta a nivel discreto,
esencial para la captura de choques. El vector fuente contiene únicamente el
acople resistivo de Ohm, los sumideros telegráficos y los términos de
Godunov--Powell, que preservan la compatibilidad termodinámica frente a los
monopolos transitorios. El jacobiano del sistema, de catorce por catorce,
posee ocho autovalores degenerados a más y menos la velocidad de la luz; su
espectro completo está en el material de respaldo.
\end{parlamento}

\lamina{10}{Mapa de Límites: del RRMHD a la Física Conocida}{Sebastián Rodriguez}{1:30}{15:30}

\begin{parlamento}{Sebastián Rodriguez}
Este mapa sitúa nuestro sistema en la física conocida. Tomando el límite de
conductividad infinita se recupera la magnetohidrodinámica ideal
relativista; con velocidades pequeñas, la magnetohidrodinámica clásica; sin
campo, la hidrodinámica de Euler; y en el límite de conductividad nula, las
ecuaciones de Maxwell en vacío con la hidrodinámica desacoplada. La misma
escalera existe para la inestabilidad que nos ocupa: en hidrodinámica la
discontinuidad es incondicionalmente inestable; en magnetohidrodinámica
clásica la tensión estabiliza si la componente del campo paralela al vector
de onda es suficiente; y en el régimen relativista ideal rige el criterio de
Mach relativista entre cero y raíz de dos.

El punto que define este trabajo es el peldaño que falta: en
magnetohidrodinámica relativista resistiva no existe relación de dispersión
algebraica cerrada para la discontinuidad tangencial, porque el término
difusivo es parabólico y suaviza la interfaz en tiempo arbitrariamente
corto. Por consiguiente, la tasa de crecimiento en función de la
conductividad debe caracterizarse numéricamente; la teoría lineal ideal
queda como referencia del límite asintótico. Esta es la justificación
epistemológica del enfoque del trabajo.
\end{parlamento}

\begin{indicaciones}
Recorrer el diagrama de izquierda a derecha con el apuntador. La lámina
también contiene la respuesta al «límite no relativista del sistema»
(Anexo A.5).
\end{indicaciones}

\lamina{11}{El Observable: Vorticidad y Enstrofía de Perturbación}{Sebastián Rodriguez}{1:30}{17:00}

\begin{parlamento}{Sebastián Rodriguez}
Antes de mostrar resultados, definimos con precisión qué se mide. El
observable primario es la enstrofía de perturbación: a la vorticidad se le
sustrae el promedio de la configuración inicial, para no contaminar la señal
con la cizalla base, y se integra su cuadrado en el dominio. En la fase
lineal, cada modo crece exponencialmente y, como la enstrofía es cuadrática
en la amplitud, el logaritmo de la enstrofía crece con pendiente igual a dos
veces la tasa de crecimiento; la regresión se efectúa en la ventana canónica
de tiempo entre dos punto cuatro y tres punto cuatro. Este factor dos,
explícito, evita cualquier ambigüedad en la extracción.

Complementan este estimador la enstrofía total y su fracción fuera del
plano, que diagnostican estructuras secundarias; el máximo de la densidad de
corriente, que mide el adelgazamiento de las láminas; la integral del
producto del campo eléctrico por la corriente, que es la disipación óhmica
global; y las energías magnética, cinética e interna. Para comparar con la
literatura, la tasa se adimensionaliza con el espesor de la capa y la
velocidad de cizalla. Ningún número que verán en adelante carece de una
definición operativa dada en esta lámina. Con el sistema y el observable
definidos, devuelvo la palabra a Bryan Martinez para la implementación
numérica.
\end{parlamento}

\relevo{Fin del segmento II. Minuto de control: 17:00. Bryan Martinez
retoma directamente con la lámina 12.}

%==============================================================================
\section{Segmento III --- Implementación numérica y montaje (Bryan Martinez, 7:30)}
%==============================================================================

\lamina{12}{Volúmenes Finitos: la Formulación Integral Obligatoria}{Bryan Martinez}{1:30}{18:30}

\begin{parlamento}{Bryan Martinez}
Gracias, Sebastián. ¿Por qué volúmenes finitos? Porque en flujos relativistas
emergen choques que vuelven singular la forma diferencial de las ecuaciones:
la formulación integral es matemáticamente obligatoria, pues admite
soluciones débiles, y los flujos numéricos respetan las condiciones de salto
de Rankine--Hugoniot. En esta formulación, la evolución del promedio de cada
celda es el balance de los flujos de Riemann en sus caras, de modo que la
conservación es exacta por construcción: el flujo que sale de una celda
entra íntegramente en la vecina. Cabe la comparación con elementos finitos:
la formulación débil global es óptima para dominios complejos con soluciones
suaves; para sistemas hiperbólicos de leyes de conservación con
discontinuidades, la captura conservativa de choques exige volúmenes
finitos.
\end{parlamento}

\lamina{13}{Reconstrucción MP5 y Solucionador HLL}{Bryan Martinez}{1:00}{19:30}

\begin{parlamento}{Bryan Martinez}
Dentro de cada celda, la reconstrucción de los estados es de quinto orden,
del tipo MP5 de Suresh y Huynh: los esquemas TVD de segundo orden recortan
los extremos y, con ello, destruyen precisamente la enstrofía que queremos
medir; MP5 la retiene de manera casi espectral. Los flujos de interfaz se
calculan con el solucionador HLL, robusto, positivo y libre de jacobianos;
su difusión sobre la onda de contacto queda compensada por la agudeza de los
estados de quinto orden, como demostró el propio grupo del director. El
criterio de diseño del estudio se enuncia en una línea: la disipación
numérica debe ser menor que la disipación física en todo el barrido; la
evidencia empírica de que así ocurre se presentará con los resultados.
\end{parlamento}

\begin{indicaciones}
Comprimible a 30 segundos si hay retraso. La pregunta previsible «¿por qué
no HLLC?» tiene respuesta dedicada en el respaldo B7 y en el Anexo A.6; no
se aborda aquí salvo que el jurado la formule.
\end{indicaciones}

\lamina{14}{Rigidez Resistiva e Integración IMEX--RK}{Bryan Martinez}{1:30}{21:00}

\begin{parlamento}{Bryan Martinez}
Como anticipó mi compañero, el término de Ohm vuelve el sistema rígido: el
tiempo de relajación resistivo es inversamente proporcional a la
conductividad y, a conductividades altas, resulta muchísimo menor que el
paso temporal hidrodinámico. Un esquema totalmente explícito exigiría pasos
de tiempo proporcionales al cuadrado del espaciamiento dividido por la
conductividad: la parálisis numérica. La solución es la partición
implícito-explícita de Runge--Kutta, en la variante fuertemente estable de
Pareschi y Russo: los flujos hiperbólicos se integran explícitamente, y la
fuente rígida, implícitamente. Y aquí está el punto técnico clave: como la
fuente no acopla celdas vecinas, la actualización implícita se reduce a un
sistema algebraico local de tres por tres con solución analítica; no hay
método de Newton global ni inversión matricial iterativa, de modo que el
sobrecosto es despreciable. Además, el esquema posee preservación
asintótica: al tomar el límite de conductividad infinita en la actualización
implícita se recupera exactamente la ley de Ohm ideal. El esquema respeta el
límite físico por construcción.
\end{parlamento}

\lamina{15}{Recuperación de Variables Primitivas: Ferrari--Cardano}{Bryan Martinez}{1:00}{22:00}

\begin{parlamento}{Bryan Martinez}
Queda un eslabón: los flujos requieren las variables primitivas, pero el
factor de Lorentz acopla cinemática y termodinámica y el mapeo conservadas a
primitivas carece de inversa explícita. Cueva reduce el problema a una
cuártica en el factor de Lorentz, que se resuelve analíticamente mediante la
transformación de Tschirnhaus y el método de Ferrari--Cardano, con un
refinamiento final de Newton--Raphson y guardias físicas: rechazo de
velocidades superlumínicas y de presiones no positivas, y una rama
linealizada estable a bajas velocidades. Esta rutina se ejecuta en cada
celda y en cada subpaso; su robustez frente a láminas de corriente agudas es
la que define el límite práctico del barrido y explica las exclusiones que
declararé enseguida.
\end{parlamento}

\begin{indicaciones}
Comprimible a 30 segundos: basta la frase «la inversión se reduce a una
cuártica resuelta analíticamente con guardias físicas», más la conexión con
las exclusiones.
\end{indicaciones}

\lamina{16}{Montaje Experimental: Doble Capa de Cizalla tipo Mizuno}{Bryan Martinez}{1:30}{23:30}

\begin{parlamento}{Bryan Martinez}
El montaje es la doble capa de cizalla de Mizuno. El perfil base es una
tangente hiperbólica doble centrada en más y menos cero coma cinco, con
velocidad de cizalla de la mitad de la velocidad de la luz y espesor de cero
coma cero cinco. Subrayo un punto importante: la semilla de la inestabilidad
es una perturbación senoidal en la velocidad transversal, el modo
fundamental compatible con la caja, con confinamiento gaussiano; el
contraste de densidad ---un flujo central diluido en un ambiente diez veces
más denso--- modela la configuración chorro--medio, pero no gatilla la
inestabilidad. El campo magnético inicial tiene una componente de guía
dominante perpendicular al plano: aporta presión pero no tensión, una
decisión topológica que permite que la inestabilidad se desarrolle y aísla
el efecto resistivo. La malla es de quinientos doce por doscientos
cincuenta y seis celdas, con unas trece celdas por capa de cizalla,
condiciones periódicas en la dirección del flujo y abiertas en la
transversal. Para dar escala astrofísica: con una longitud de referencia de
diez a la dieciséis centímetros, el tiempo característico es de unos cuatro
días, y el vórtice maduro alcanza unas ciento setenta unidades
astronómicas.
\end{parlamento}

\begin{indicaciones}
Señalar en la figura las dos interfaces rectas en el estado inicial. La
precisión «la densidad no gatilla nada» corrige una lectura errónea
previsible del montaje.
\end{indicaciones}

\lamina{17}{Matriz de Simulaciones y Criterios de Exclusión}{Bryan Martinez}{1:00}{24:30}

\begin{parlamento}{Bryan Martinez}
Esta es la matriz completa de simulaciones: el barrido de cuarenta y cuatro
valores de conductividad con número de Courant de cero coma uno, y las
campañas magnéticas a conductividad seis mil y diez mil, con Courant
reducido a cero coma cero cuatro y cero coma cero dos, porque las láminas de
corriente agudas exigen mayor estabilidad en la recuperación de primitivas.
En total, unas treinta mil horas-núcleo de cómputo. Y declaramos
explícitamente los datos excluidos: en la campaña A, los dos casos de campo
más débil fallaron por inestabilidad numérica antes de la fase no lineal; en
la campaña B, el ángulo de quince grados se detuvo por un fallo de
convergencia del solucionador de primitivas; y en el barrido, el análisis
cuantitativo de la tasa lineal usa las simulaciones con conductividad mayor
o igual a mil seiscientos, con coeficiente de determinación de al menos cero
coma noventa. Preferimos declarar estas exclusiones a ocultarlas: son parte
del resultado. Con esto, mi compañero presenta los resultados del barrido en
conductividad.
\end{parlamento}

\relevo{Fin del segmento III. Minuto de control: 24:30. Sebastián Rodriguez
retoma con la lámina 18.}

%==============================================================================
\section{Segmento IV --- Resultados I: barrido en conductividad (Sebastián Rodriguez, 13:00)}
%==============================================================================

\lamina{18}{El Barrido Completo en Conductividad}{Sebastián Rodriguez}{1:00}{25:30}

\begin{parlamento}{Sebastián Rodriguez}
Gracias, Bryan. Esta es la vista panorámica del barrido: la enstrofía de
perturbación en el tiempo para las cuarenta y cuatro simulaciones, con la
conductividad codificada en color de negro a naranja. Obsérvese la
estratificación monótona y limpia de las curvas. Sobre este conjunto se
definen los subconjuntos de análisis, con criterios explícitos: el ajuste de
la tasa lineal emplea treinta y cinco simulaciones, aquellas con
conductividad mayor o igual a mil seiscientos y ajuste de calidad; los
ajustes de pico emplean veintinueve, porque el pico solo es identificable
dentro del horizonte temporal a partir de conductividad dos mil
ochocientos. La pendiente lineal es medible antes que el pico; de ahí la
diferencia entre ambos subconjuntos.
\end{parlamento}

\lamina{19}{Visión Global del Barrido: Fenomenología}{Sebastián Rodriguez}{1:00}{26:30}

\begin{parlamento}{Sebastián Rodriguez}
Cualitativamente, el barrido exhibe tres regímenes temporales en todos los
casos: fase lineal hacia tiempo tres, enrollamiento no lineal hacia tiempo
ocho y saturación hacia tiempo catorce. El contraste morfológico entre
columnas es el mensaje: a conductividad cien, la difusión óhmica amortigua y
difumina los vórtices; a conductividad diez mil, con el campo casi
congelado, aparecen subestructura fina y turbulencia de pequeña escala. Las
cuarenta y cuatro simulaciones alcanzaron tiempos superiores a cinco sin
inestabilidades numéricas espurias.
\end{parlamento}

\begin{indicaciones}
Comprimible a 30 segundos: basta recorrer la figura señalando las tres
columnas. No detenerse en detalles de paneles.
\end{indicaciones}

\lamina{20}{Extracción de la Tasa de Crecimiento Lineal}{Sebastián Rodriguez}{1:30}{28:00}

\begin{parlamento}{Sebastián Rodriguez}
La extracción cuantitativa procede como se definió: regresión del logaritmo
de la enstrofía en la ventana canónica, con pendiente igual a dos veces la
tasa. Los criterios de calidad son explícitos: la ventana es fija,
posterior al transitorio inicial y anterior a la saturación; el coeficiente
de determinación clasifica los regímenes en resistivo, de transición e
ideal; y el análisis cuantitativo emplea las treinta y cinco simulaciones
con conductividad de mil seiscientos en adelante. Una precisión conceptual
importante: por debajo de conductividad alrededor de mil cuatrocientos no
existe régimen lineal limpio; no se trata de un ajuste pobre, sino de que
la difusión domina desde el instante inicial y el modo no alcanza a
establecerse. Reconocer ese límite es parte del rigor del análisis, no una
debilidad.
\end{parlamento}

\lamina{21}{La Transición Resistivo--Ideal: Sigmoide con Offset}{Sebastián Rodriguez}{2:00}{30:00}

\begin{parlamento}{Sebastián Rodriguez}
Llegamos al resultado paramétrico central del objetivo uno. La dependencia
de la tasa de crecimiento con la conductividad queda descrita por un modelo
sigmoide con desplazamiento, de cuatro parámetros. Dos números resumen el
ajuste: la asíntota ideal, suma del desplazamiento y de la amplitud, vale
uno coma cero tres, con incertidumbre de cero coma cero cinco; y el centro
de la transición se sitúa en conductividad dos mil ochocientos quince, con
error de ciento veintitrés.

La validación se sostiene en tres argumentos independientes. Primero, el
criterio de información de Akaike en comparación anidada: agregar el
parámetro de desplazamiento, sobre los mismos datos, mejora el criterio en
ciento veintiséis unidades, evidencia estadísticamente decisiva. Segundo,
validación fuera de muestra: los nueve puntos resistivos que no
participaron del ajuste se predicen un treinta y seis por ciento mejor con
el modelo completo. Tercero, honestidad sobre los residuos: están
autocorrelacionados, de modo que el error formal subestima la incertidumbre
real; precisamente por eso la conductividad crítica se estima con múltiples
métodos, como muestro a continuación. Subrayo, además, que el
desplazamiento es un parámetro del modelo con soporte predictivo, no una
tasa de crecimiento negativa con interpretación física directa.
\end{parlamento}

\begin{indicaciones}
Lámina central: no acelerar. Los tres argumentos de validación se enumeran
con los dedos. El detalle estadístico completo está en el respaldo B8.
\end{indicaciones}

\lamina{22}{Los Cuatro Estimadores de la Conductividad Crítica}{Sebastián Rodriguez}{1:00}{31:00}

\begin{parlamento}{Sebastián Rodriguez}
¿Dónde ocurre la transición? Lo preguntamos a cuatro observables
independientes, con un método común: el punto de inflexión de cada curva
respecto del logaritmo de la conductividad, con errores por remuestreo. El
sigmoide de la tasa lineal responde dos mil ochocientos quince; el máximo
de enstrofía, tres mil cuatrocientos; el tiempo del pico, seis mil; y la
tasa en escala logarítmica, seis mil ochocientos. Cuatro respuestas
distintas y legítimas: esa dispersión es el dato que motiva la lámina
siguiente.
\end{parlamento}

\begin{indicaciones}
Comprimible a 30 segundos. Los valores están escritos en pantalla; no es
necesario repetir los errores individuales de viva voz.
\end{indicaciones}

\lamina{23}{La Conductividad Crítica No Es un Número: Es una Zona}{Sebastián Rodriguez}{1:30}{32:30}

\begin{parlamento}{Sebastián Rodriguez}
La lectura conjunta es esta cascada de activación: cada observable se
idealiza a una conductividad distinta, desde el máximo de estructuras fuera
del plano alrededor de mil cuatrocientos hasta la saturación no lineal
cerca de siete mil. La conclusión honesta es que la conductividad crítica
no es un número, sino una zona de transición: el intervalo de mil
cuatrocientos a siete mil, con centro en torno a dos mil ochocientos
quince, que en número de Lundquist corresponde a un orden de mil. Y el
mensaje epistemológico es explícito: la dispersión entre métodos, de más o
menos mil novecientos, no es ruido; es la firma de una transición
multiescala, en la que cada fenómeno físico se congela a su propia escala.
Reportar la banda completa ---y no un falso valor único con un error formal
de ciento veintitrés--- es lo científicamente riguroso.
\end{parlamento}

\lamina{24}{Contraste con la Teoría Lineal: la Escalera de Predicciones}{Sebastián Rodriguez}{2:00}{34:30}

\begin{parlamento}{Sebastián Rodriguez}
Este es, a nuestro juicio, el resultado central del trabajo. ¿Es correcto el
valor asintótico medido? La cadena lógica tiene tres peldaños. Primero, el
techo hidrodinámico inviscido: la ecuación de Rayleigh para el perfil de
tangente hiperbólica, con nuestro solucionador validado contra el resultado
clásico de Michalke, da una tasa normalizada máxima de cero coma ciento
noventa. Segundo, la geometría fija la física: como el vector de onda es
perpendicular al campo de guía, el campo aporta presión pero no tensión, y
la velocidad de restitución no es la del sonido sino la magnetosónica
rápida perpendicular, cero coma seiscientos noventa y uno; con ella, la
teoría lineal compresible de Lees y Lin predice cero coma cero noventa y
cinco. Tercero, nuestro valor medido: cero coma ciento tres. El acuerdo es
del ocho por ciento, sin ningún parámetro de ajuste. El control negativo
refuerza la cadena: si se usara la velocidad del sonido sola, la predicción
colapsaría a cero coma cero dieciséis; es la presión magnética del campo de
guía la que sostiene la inestabilidad. Además, el número de Mach
relativista vale cero coma sesenta, holgadamente por debajo del umbral de
estabilización de raíz de dos.

Declaramos el alcance con precisión: se trata de una consistencia
cuantitativa fuerte, no de una validación cerrada, porque la asíntota es
una extrapolación ---nuestro dato máximo alcanza el ochenta y cinco por
ciento del plateau--- y la predicción emplea la sustitución escalar de la
velocidad característica, no la relación de dispersión completa de grado
ocho.
\end{parlamento}

\begin{indicaciones}
Lámina de mayor peso: dicción pausada en los tres peldaños de la escalera,
señalándolos sobre la figura. La frase final de alcance se pronuncia
literalmente; es el blindaje de honestidad del resultado.
\end{indicaciones}

\lamina{25}{No Linealidad: Ley de Potencia de la Enstrofía Máxima}{Sebastián Rodriguez}{1:30}{36:00}

\begin{parlamento}{Sebastián Rodriguez}
En el régimen no lineal, objetivo dos, el máximo de la enstrofía de
perturbación sigue una ley de potencia en la conductividad con exponente
uno coma veintidós y coeficiente de determinación de cero coma novecientos
noventa y siete. El punto de referencia teórico es el piso de una lámina de
corriente única de Sweet--Parker, cuyo espesor escala como la raíz del
número de Lundquist y daría exponente un medio; ese escalamiento es robusto
frente a correcciones relativistas. Nuestro exponente lo excede en un
factor de dos coma cuatro. La interpretación se declara expresamente como
hipótesis fenomenológica: el exceso sugiere que la lámina primaria se
fragmenta por el modo de desgarro ---el \emph{tearing}--- multiplicando la
superficie disipativa disponible; la firma es consistente con la
literatura, pero su demostración cuantitativa exigiría resolver la
jerarquía de sub-láminas con esquemas de orden superior, y queda planteada
como trabajo futuro.
\end{parlamento}

\lamina{26}{Variables Globales: la Firma Electromagnética de la Conductividad}{Sebastián Rodriguez}{1:30}{37:30}

\begin{parlamento}{Sebastián Rodriguez}
El objetivo tres se responde con los diagnósticos globales. El trabajo
electromagnético disipado acumulado crece de cero coma cero tres a cero
coma cuarenta y cuatro a lo largo del barrido: a mayor conductividad, más
violencia hidrodinámica procesa más energía magnética. La amplificación de
la energía magnética propia decae del quince por ciento a cero hacia el
límite ideal: el congelamiento impide reorganizar la topología. Y la
fracción de estructuras fuera del plano alcanza su máximo, cero coma
treinta y siete, justo en la zona de transición.

Y aquí, señor jurado, está la evidencia empírica que anunciamos: el panel
de la densidad de corriente máxima exhibe un plateau que marca dónde la
difusividad física supera el error de truncamiento de la malla. Esto
demuestra que los resultados en la zona de transición no están contaminados
por disipación artificial. Más aún: se sabe, por Lecoanet y colaboradores,
que la inestabilidad de Kelvin--Helmholtz bidimensional ideal no converge
con la resolución; es precisamente la resistividad explícita la que da
significado convergente a la tasa de crecimiento como función de la
conductividad. La disipación de nuestras simulaciones es física, no
numérica. Con los resultados del barrido completos, Bryan Martinez
presenta la campaña magnética y las conclusiones.
\end{parlamento}

\relevo{Fin del segmento IV. Minuto de control: 37:30. Bryan Martinez
retoma con la lámina 27.}

%==============================================================================
\section{Segmento V --- Resultados II: campaña magnética y cierre (Bryan Martinez, 7:00)}
%==============================================================================

\lamina{27}{Campaña Magnética I: Intensidad y Orientación}{Bryan Martinez}{1:30}{39:00}

\begin{parlamento}{Bryan Martinez}
Gracias, Sebastián. El objetivo cuatro pregunta quién suprime la inestabilidad:
¿la presión magnética o la tensión? La campaña A, de intensidad, responde
lo primero: al duplicar el campo, la tasa de crecimiento cae de cero coma
ochenta y uno a cero coma treinta y ocho, y la enstrofía máxima, de
ochenta y siete a veinticinco. Pero el flujo permanece súper-alfvénico en
la dirección paralela en todos los casos: la tensión es incapaz de frenar
la cizalla, de modo que el agente supresor es la presión magnética, es
decir, la magnetización creciente. La campaña B, de orientación en el
plano que contiene al vector de onda, muestra que la tensión solo actúa a
través de la componente paralela a dicho vector. Destaco la aparente
paradoja del ángulo cero: campo de guía puro, sin tensión, y sin embargo
enstrofía mínima; la resolución es que, sin componente en el plano, el
vórtice bidimensional no puede estirar líneas de campo, y sin estiramiento
no hay dínamo ni reconexión que realimenten la dinámica. Y en el ángulo de
noventa grados, el máximo local proviene de capas de reconexión finas
mientras la enstrofía integrada sí decrece con la tensión: un efecto del
observable, no de la física.
\end{parlamento}

\lamina{28}{Campaña Magnética II: Balance Energético y Canales de Disipación}{Bryan Martinez}{1:30}{40:30}

\begin{parlamento}{Bryan Martinez}
Cuando la geometría enmascara la fase lineal, el observable correcto no es
la tasa de crecimiento sino la energía. Tres hallazgos. Primero, las
energías cinética y magnética oscilan en anti-fase, ciclo a ciclo, con
correlaciones entre menos cero coma seis y menos cero coma nueve: existe
un canal de conversión directo y reversible, el dínamo de ida y el rebote
elástico de la tensión de vuelta. Segundo, la resistividad decide el
destino final de esa energía: a conductividad diez mil el intercambio es
cuasi-reversible; a seis mil, la disipación óhmica desvía energía a calor,
hasta un orden de magnitud más en la orientación perpendicular. Tercero,
como control de calidad interno, la energía total se conserva con error
relativo menor que una parte en mil en todas las simulaciones: una
validación independiente del esquema numérico.
\end{parlamento}

\lamina{29}{Campaña Magnética III: Orientación sin Tensión}{Bryan Martinez}{1:00}{41:30}

\begin{parlamento}{Bryan Martinez}
La campaña C es el experimento de control: al anular la componente del
campo en la dirección del flujo, el producto del vector de onda con el
campo es nulo y la tensión desaparece, mientras el módulo del campo ---y
con él la presión magnética--- permanece constante con el ángulo. El
resultado es contundente: sin tensión no hay fase lineal limpia; el
crecimiento pierde su carácter exponencial y la mezcla se vuelve difusa,
de manera idéntica en ambas conductividades. La tensión magnética es, por
tanto, el estabilizador primario de la inestabilidad en este régimen. Y
con la tensión ausente, la conversión energética queda al desnudo: la
anti-correlación entre las energías cinética y magnética se intensifica
hasta menos cero coma noventa y ocho en los ángulos altos ---conversión
casi uno a uno--- con disipación óhmica monótona en el ángulo.
\end{parlamento}

\lamina{30}{Conclusiones (por objetivo)}{Bryan Martinez}{1:30}{43:00}

\begin{parlamento}{Bryan Martinez}
Concluimos en correspondencia con los cuatro objetivos. Sobre el primero:
la transición del régimen resistivo al ideal queda descrita por el modelo
sigmoide con desplazamiento; la asíntota ideal es uno coma cero tres más
menos cero coma cero cinco, consistente al ocho por ciento con la teoría
lineal compresible, y la conductividad crítica es una zona, de mil
cuatrocientos a siete mil, con centro cerca de dos mil ochocientos.
Sobre el segundo: la enstrofía máxima escala con exponente uno coma
veintidós, excede el piso de Sweet--Parker y sugiere desgarro secundario
de las láminas; las estructuras tridimensionales emergen donde difusión y
advección compiten. Sobre el tercero: la conductividad controla la
intensidad de las láminas de corriente, la disipación acumulada y la
amplificación de la energía magnética propia. Y sobre el cuarto: la
supresión por intensidad es un efecto de presión y magnetización; la
tensión actúa únicamente a través de la componente del campo paralela al
vector de onda; sin tensión dominan la reconexión y la disipación
directa; y la anti-fase entre energías cinética y magnética constituye el
canal universal de conversión, cuyo destino ---intercambio reversible o
calor--- lo decide la resistividad.
\end{parlamento}

\begin{indicaciones}
Se enuncian los resultados sin autoevaluación: la valoración del
cumplimiento de los objetivos corresponde al jurado.
\end{indicaciones}

\lamina{31}{Limitaciones del Modelo y Trabajo Futuro}{Bryan Martinez}{1:00}{44:00}

\begin{parlamento}{Bryan Martinez}
Declaramos los límites del trabajo, que son a la vez su mapa de
continuación. El estudio es bidimensional, con ecuación de estado
politrópica; la conductividad es escalar, sin efecto Hall ni anisotropía
---el modelo de fluido único trunca la jerarquía de momentos de Vlasov y
es válido mientras la escala disipativa relevante sea la capa resistiva y
no el girorradio---; a las conductividades más altas las láminas de
corriente rozan el límite de resolución de la malla; y la asíntota ideal
es una extrapolación, pues el dato máximo alcanza el ochenta y cinco por
ciento del plateau. En consecuencia, el trabajo futuro contempla: medir
directamente el plateau con simulaciones de mayor conductividad; resolver
la relación de dispersión completa de grado ocho en estos parámetros;
cartografiar con más ángulos la frontera entre tensión y presión; y
extender el estudio a tres dimensiones con esquemas de cuarto orden,
capaces de resolver las hojas resistivas.
\end{parlamento}

\lamina{32}{Agradecimiento y apertura del turno de preguntas}{Bryan Martinez}{0:30}{44:30}

\begin{parlamento}{Bryan Martinez}
Señor jurado, señor director, señoras y señores: agradecemos su atención.
Quedamos, con mi compañero, a su entera disposición para las preguntas y
la discusión que consideren pertinentes.
\end{parlamento}

\begin{indicaciones}
Ambos ponentes de pie durante el turno de preguntas. La lámina de
referencias no se expone: se salta directamente a la de agradecimiento. El
material de respaldo B1--B13 queda listo para navegar según la pregunta
(mapa en el Anexo B).
\end{indicaciones}

%==============================================================================
\appendix
\section{Anexo A --- Respuestas ensayadas de sesenta segundos}
%==============================================================================

Se transcriben las seis respuestas preparadas para el turno de preguntas,
derivadas del perfil del jurado. Cada una está calibrada a un minuto y
puede ampliarse con la lámina de respaldo indicada.

\subsection{«Explique la inestabilidad sin fórmulas»}
Dos flujos en contacto a velocidades distintas almacenan energía libre en
su capa de cizalla. Una ondulación mínima de la interfaz acelera el flujo
sobre las crestas y lo frena en los valles; por Bernoulli, esa asimetría
de presiones ejerce un torque que amplifica la ondulación. La
retroalimentación produce crecimiento exponencial y, al saturar, la
interfaz se enrolla en vórtices que mezclan ambos medios. En un chorro
astrofísico, ese mecanismo frena el chorro, mezcla materia del ambiente y
amplifica el campo magnético. \emph{(Lámina 3.)}

\subsection{«¿Cómo garantiza la ausencia de monopolos sin destruir la causalidad?»}
En el continuo no hay nada que garantizar: la divergencia nula es una
identidad geométrica, consecuencia de la nilpotencia de la derivada
exterior. La malla discreta rompe esa nilpotencia y genera monopolos
espurios. El método GLM acopla las restricciones a la evolución mediante
dos escalares con origen variacional; el error resultante obedece la
ecuación del telégrafo: se propaga hiperbólicamente a la velocidad de
limpieza ---dentro del cono causal--- y decae exponencialmente. La
causalidad queda intacta por construcción. \emph{(Lámina 8; respaldo B3.)}

\subsection{«¿Su disipación es física o numérica?»}
Tres argumentos. De diseño: reconstrucción MP5 de quinto orden con flujo
HLL, elegida para que la disipación numérica sea menor que la física en
todo el barrido. Empírico: el plateau de la densidad de corriente máxima
marca dónde la difusividad física supera el error de truncamiento; los
resultados de la zona de transición están dentro de ese régimen.
Conceptual: la KHI bidimensional ideal ni siquiera converge con la
resolución (Lecoanet 2016); es la resistividad explícita la que da
significado convergente a la tasa medida. \emph{(Láminas 13 y 26;
respaldo B7.)}

\subsection{«¿De dónde sale la resistividad? ¿Cuál es su origen microscópico?»}
La resistividad es un coeficiente de transporte fenomenológico que emerge
de las integrales de colisión del nivel cinético de Vlasov--Boltzmann al
truncar la jerarquía de momentos en el fluido único. El cierre escalar e
isótropo exige que la frecuencia de colisiones domine sobre la frecuencia
de ciclotrón; el modelo es válido mientras la escala disipativa relevante
sea la capa resistiva, del orden del inverso de la raíz del número de
Lundquist, y no el girorradio. Fuera de ese régimen se requerirían efecto
Hall, anisotropía o descripción cinética. \emph{(Lámina 7; respaldos B5 y
B9.)}

\subsection{«¿Cómo se recupera la física conocida desde su sistema?»}
Por límites sucesivos: con conductividad infinita, el campo eléctrico
comóvil se anula y se recupera la magnetohidrodinámica ideal relativista;
con velocidades pequeñas y factor de Lorentz unidad, la
magnetohidrodinámica clásica; sin campo magnético, la hidrodinámica de
Euler; y con conductividad nula, Maxwell en vacío con la hidrodinámica
desacoplada. El integrador respeta el primer límite por construcción:
posee preservación asintótica y converge exactamente a la ley de Ohm ideal
cuando la conductividad diverge. \emph{(Láminas 10 y 14; respaldo B4.)}

\subsection{«El artículo del código usa HLLC; ¿por qué ustedes HLL?»}
Es una decisión de diseño documentada. HLLC restaura la onda de contacto,
pero HLL con reconstrucción MP5 de quinto orden logra un efecto
equivalente por otra vía: la agudeza de los estados reconstruidos compensa
la difusión del solucionador sobre el contacto ---hallazgo del propio
grupo del director--- con mayor simplicidad, positividad garantizada y
menor costo en un barrido de sesenta y dos simulaciones. La verificación
empírica de que la disipación numérica quedó por debajo de la física es el
plateau de la corriente máxima. \emph{(Lámina 13; respaldo B7.)}

%==============================================================================
\section{Anexo B --- Mapa del material de respaldo}
%==============================================================================

\begin{center}
\begin{tabular}{@{}cll@{}}
\toprule
Respaldo & Contenido & Se invoca si preguntan por\ldots\\
\midrule
B1  & Tensor de Faraday desde el cuadripotencial & formalismo tensorial, Bianchi\\
B2  & Tensor energía-momento por vía de Hilbert  & origen variacional de $T^{\mu\nu}$\\
B3  & GLM desde la acción extendida              & telégrafo, forma de Dedner\\
B4  & Espectro característico $14\times14$       & autovalores, CFL, estabilidad\\
B5  & Proyección $3+1$ de la ley de Ohm          & álgebra de la ley de Ohm\\
B6  & Validación del solucionador lineal          & Michalke, umbral $\sqrt2$\\
B7  & ¿Por qué HLL$+$MP5 y no HLLC?              & elección del solucionador\\
B8  & Estadística del ajuste sigmoide             & AIC/BIC, degeneración de parámetros\\
B9  & De Vlasov al fluido                         & fundamento cinético del continuo\\
B10 & Recursos computacionales                    & horas-núcleo, flujo de trabajo\\
B11 & Estructuras secundarias y campaña C         & repreguntas del objetivo 4\\
B12 & Mapas de densidad de las campañas           & morfología de las campañas\\
B13 & Tablero dinámico de la campaña B            & evolución comparada $\sigma=6000$/$10000$\\
\bottomrule
\end{tabular}
\end{center}

\vspace{6pt}
La lámina de referencias contiene las cuarenta y cuatro fuentes de la
monografía; no se expone, pero permanece disponible si el jurado solicita
una cita en particular.

\end{document}
